% Laget automatisk av elevark/lag_elevark.py fra opplegg/lengdekontraksjon.html. Endringer her blir overskrevet.
\documentclass[11pt]{article}
\usepackage{elevark}
\begin{document}
\elevarkhode{Fysikk 2 · Spesiell relativitetsteori}{Lengdekontraksjon – arbeidsark}{https://lektorbo.no/sim/tog-og-tunnel.html}{Tog og tunnel}

\section*{Oppgaver}

Oppgave 1–6 bruker Tog og tunnel\footnote{\url{lektorbo.no/sim/tog-og-tunnel.html}}, og oppgave 7–8 bruker Romreise\footnote{\url{lektorbo.no/sim/romreise.html}}.

\begin{oppgaveliste}
\oppgave \fase{test}{Regn}Sett farten til $0{,}80c$. Regn ut $γ$ og lengden av toget i tunnelens system. Passer toget inn i tunnelen? Kjør simuleringen og sjekk.
\svar
\oppgave \fase{gjett}{Gjett}Bytt til togets system. Hvor lang er tunnelen nå? Kan hele toget være inne samtidig? Hva tror du skjer med portene?
\svar
\oppgave \fase{forklar}{Forklar}Se på tabellen over hendelser. I hvilken rekkefølge lukkes og åpnes portene i togets system? Hvorfor blir ikke toget truffet av noen port?
\svar
\oppgave \fase{test}{Undersøk}Finn den laveste farten der toget passer inn sett fra tunnelen. Regn først, og sjekk så. Prøv også $0{,}60c$. Hva skjer da, og skjer det samme i begge systemene?
\svar
\oppgave \fase{test}{Mål}Velg «Tunnelfolket måler toget» med $0{,}80c$. Hvor langt fra hverandre står merkene? Bytt til togets system. Hvorfor mener togfolket at målingen er feil?
\svar
\oppgave \fase{test}{Mål}Velg «Togfolket måler tunnelen». Hvor langt er det mellom merkene på toget? Hvem måler den andre som kortest?
\svar
\oppgave \fase{test}{Regn}Romreise: Velg Alfa Centauri og $0{,}90c$. Regn ut reisetiden på jorda og på skipet. Hvor lang er avstanden i skipets system? Hvordan forklarer mannskapet at de kommer så fort fram?
\svar
\oppgave \fase{forklar}{Forklar}Mannskapet har reist 4,37~lysår på 2,12~år. Betyr det at de har reist fortere enn lyset?
\svar
\end{oppgaveliste}

\end{document}
